\documentclass[11pt]{report}
\usepackage[utf8]{inputenc}

% include images -- % 
\usepackage{graphicx}

% define page geometry -- %
\usepackage{geometry}

\geometry{a4paper, margin=100px}

% enables referencing to enum items -- %
\usepackage{enumitem}

% places footnotes at the bottom -- %
\usepackage[bottom]{footmisc}

% hyperlinks and hyper references -- %
\usepackage{hyperref}
\usepackage{colortbl}

\usepackage{array}
\usepackage{float}

% tables -- % 
\usepackage{tabularx}
\usepackage{longtable} % splits a table when it takes more then one page

% including pictures in the title -- %
\usepackage{titlepic}

% fancy headers -- %
\usepackage{fancyhdr}

\pagestyle{fancy}

\fancyhead[L]{\nouppercase{\rightmark}}
\fancyhead[R]{\nouppercase{\leftmark}}
\setlength{\headheight}{13.6pt}
\addtolength{\topmargin}{-0.1pt}

%table coloring -- %
\usepackage[table]{xcolor} % Add the xcolor package for background color

%table padding -- %
\renewcommand{\arraystretch}{1.5}

% hyperlinks -- %
\hypersetup{
    colorlinks=true,
    linkcolor=black,
    filecolor=magenta,      
    urlcolor=purple,
    pdftitle={Requirement Analysis and Specification Document},}

% paths to the graphic folder relative to the tex file -- %
\graphicspath{ {./images/} }


% --------------------------------------------------------------------------------------------- %

\pagenumbering{Roman}

\titlepic{\includegraphics[width=0.5\textwidth]{logo_polimi}}
\title{CodeKataBattle - Design Document}
% https://github.com/FilippoRiva/RivaSulfaro
\author{Filippo Riva, Antonio Sulfaro}
\date{\today}

\begin{document}

\maketitle

\tableofcontents

\cleardoublepage










%  1  --------------------------------------------------------------------------------------------%

\pagenumbering{arabic}

\chapter{Introduction}

\section{Purpose}
    %TODO

\section{Scope}
    %TODO

\section{Definitions, Acronyms, Abbreviations}

Here's a list of useful definitions necessary to a comprehensive understanding of the CodeKataBattle RASD document

\subsection{Definitions}

\begin{description}
    \item[Kata Tournament :] A competition organized by educators on the CKB platform, where students
        can participate in various Kata Battles.
    \item[Kata Battle :] A competition organized by an educator associated to a Kata Tournament, where students
        compete with each other by programming solutions to code katas.
    \item[Organizer :] An educator who possesses the permission to upload KB for a specific tournament.
    This authorization is granted when the educator has been invited by the creator of the tournament.
    \item[Code Kata :] A programming problem associated with each Kata Battle uploaded by an Organizer.
        It consists in a description of the problem, a software project, build automation scripts and test cases.
    \item[Team :] A group of one or more students taking part to a Kata Battle.
    \item[Consolidation Stage :] In this phase, the Kata Battle's Organizer will manually assess and evaluate the submissions made by students.
        If no manual evaluation is required, it does not exist.
        This phase commences upon the conclusion of a KB and extends until the educator decides to finalize it.
    \item[Tournament Rank :] The position of a student in a tournament in comparison to others based on the sum of the scores
    he obtained in each Kata Battle of that Kata Tournament.
    \item[Battle Rank :] The position of a team in a battle in comparison to others based on their score.
    \item[Score :] The numerical representation, between 0 and 100, of the students' performance in a battle,
        calculated taking into consideration every evaluation factor in a proper way.
    \item[Major Event :] Events that take place in the CKB platform and need to be notified to the specific students (full list in section \ref{standard_compliance}).
\end{description}

\subsection{Acronyms}

\begin{description}
    \item[CKB :] CodeKataBattle platform.
    \item[KB :] Kata Battle.
    \item[KT :] Kata Tournament.
    \item[WPx :] x-th World Phenomena.
    \item[SPx :] x-th Shared Phenomena.
    \item[Rx :] x-th Requirement.
    \item[Dx :] x-th Domain Assumption.
    \item[Sx :] x-th Scenario.
    \item[UCx :] x-th Use Case.
\end{description}

\subsection{Abbreviations}

\begin{description}
    \item[repo :] GitHub repository.
\end{description}

\section{Revision History}

\begin{center}
\begin{tabular}{|c|c|p{9cm}|}
    \hline
    \rowcolor{gray!20}
    \textbf{Version} & \textbf{Date} & \textbf{Changes} \\
    \hline
    0.1 & 12/21/2023 & Document Structure\\
    \hline
\end{tabular}
\end{center}

\section{Reference Documents}

Reference Books : 

Hans Van Vilet - Software Engineering: Principles and Practice

\section{Document Structure}

We offer a concise overview of the document structure: 
\begin{itemize}
    \item \textbf{Chapter 1: Introduction}\\
        %TODO
    \item \textbf{Chapter 2: Architectural Design}\\
        %TODO
    \item \textbf{Chapter 3: User Interface Design}\\
        %TODO
    \item \textbf{Chapter 4: Requirements Traceability}\\
        %TODO
    \item \textbf{Chapter 5: Implementation, Integration and Test Plan}\\
        %TODO
    \item \textbf{Chapter 6: Effort Spent}\\
        %TODO
\end{itemize}

%  2  --------------------------------------------------------------------------------------------%

\chapter{Architectural Design}


\section{Overview}
\section{Component view}
\section{Deployment view}
\section{Runtime view}
\section{Component interfaces}
\section{Selected architectural styles and patterns}
\section{Other design decisions}

%  3  --------------------------------------------------------------------------------------------%

\chapter{User Interface Design}


%  4  --------------------------------------------------------------------------------------------%

\chapter{Requirements Traceability}


%  5  --------------------------------------------------------------------------------------------%

\chapter{Implementation, Integration and Test Plan}



%  6  --------------------------------------------------------------------------------------------%
\chapter{Effort Spent}
\begin{table}[H]
\centering
\begin{tabular}{|l|l|}
\hline
\textbf{Student}       & \textbf{Effort} \\ \hline
Filippo Riva           & 0h              \\ \hline
Antonio Sulfaro        & 0h              \\ \hline
\end{tabular}
\end{table}

\end{document}